\begin{animation}[id="planets-queries-2", label="CSES Planets Queries II -- Functional Graph Distance"]
\shape{g}{Graph}{nodes=[1,2,3,4,5,6], edges=[(1,2),(2,3),(3,1),(4,5),(5,6),(6,4)], directed=true, layout="stable", layout_seed=7}
\shape{lift}{DPTable}{rows=3, cols=6, data=["","","","","","","","","","","","","","","","","",""]}
\shape{ans}{Array}{size=1, data=["?"], label="Answer"}

\step
\narrate{Functional graph: each node has exactly one outgoing edge. Node 1 teleports to 2, node 2 to 3, node 3 to 1 (cycle of length 3). Similarly 4->5->6->4 forms a second cycle. Each weakly connected component is rho-shaped -- here both are pure cycles.}

\step
\recolor{g.node[1]}{state=current}
\recolor{g.node[2]}{state=current}
\recolor{g.node[3]}{state=current}
\recolor{g.edge[(1,2)]}{state=current}
\recolor{g.edge[(2,3)]}{state=current}
\recolor{g.edge[(3,1)]}{state=current}
\narrate{Component 1: cycle 1 -> 2 -> 3 -> 1 with length 3. All three nodes are on the cycle (depth 0).}

\step
\recolor{g.node[1]}{state=idle}
\recolor{g.node[2]}{state=idle}
\recolor{g.node[3]}{state=idle}
\recolor{g.edge[(1,2)]}{state=idle}
\recolor{g.edge[(2,3)]}{state=idle}
\recolor{g.edge[(3,1)]}{state=idle}
\recolor{g.node[4]}{state=current}
\recolor{g.node[5]}{state=current}
\recolor{g.node[6]}{state=current}
\recolor{g.edge[(4,5)]}{state=current}
\recolor{g.edge[(5,6)]}{state=current}
\recolor{g.edge[(6,4)]}{state=current}
\narrate{Component 2: cycle 4 -> 5 -> 6 -> 4 with length 3. Also a pure cycle. Nodes in different components can never reach each other.}

\step
\recolor{g.node[4]}{state=idle}
\recolor{g.node[5]}{state=idle}
\recolor{g.node[6]}{state=idle}
\recolor{g.edge[(4,5)]}{state=idle}
\recolor{g.edge[(5,6)]}{state=idle}
\recolor{g.edge[(6,4)]}{state=idle}
\narrate{Now we build the binary lifting table. lift[k][i] = the node reached from i after 2^k teleportations. Row 0 is k=0 (1 step), row 1 is k=1 (2 steps), row 2 is k=2 (4 steps).}

\step
\apply{lift.cell[0][0]}{value="2"}
\apply{lift.cell[0][1]}{value="3"}
\apply{lift.cell[0][2]}{value="1"}
\apply{lift.cell[0][3]}{value="5"}
\apply{lift.cell[0][4]}{value="6"}
\apply{lift.cell[0][5]}{value="4"}
\recolor{lift.cell[0][0]}{state=current}
\recolor{lift.cell[0][1]}{state=current}
\recolor{lift.cell[0][2]}{state=current}
\recolor{lift.cell[0][3]}{state=current}
\recolor{lift.cell[0][4]}{state=current}
\recolor{lift.cell[0][5]}{state=current}
\narrate{Row k=0: lift[0][i] = t[i], the direct teleport target. Node 1->2, 2->3, 3->1, 4->5, 5->6, 6->4. This is just the original edge list.}

\step
\recolor{lift.cell[0][0]}{state=done}
\recolor{lift.cell[0][1]}{state=done}
\recolor{lift.cell[0][2]}{state=done}
\recolor{lift.cell[0][3]}{state=done}
\recolor{lift.cell[0][4]}{state=done}
\recolor{lift.cell[0][5]}{state=done}
\apply{lift.cell[1][0]}{value="3"}
\apply{lift.cell[1][1]}{value="1"}
\apply{lift.cell[1][2]}{value="2"}
\apply{lift.cell[1][3]}{value="6"}
\apply{lift.cell[1][4]}{value="4"}
\apply{lift.cell[1][5]}{value="5"}
\recolor{lift.cell[1][0]}{state=current}
\recolor{lift.cell[1][1]}{state=current}
\recolor{lift.cell[1][2]}{state=current}
\recolor{lift.cell[1][3]}{state=current}
\recolor{lift.cell[1][4]}{state=current}
\recolor{lift.cell[1][5]}{state=current}
\narrate{Row k=1: lift[1][i] = lift[0][lift[0][i]], the node reached after 2 steps. For node 1: lift[0][lift[0][1]] = lift[0][2] = 3. Each entry chains two single-step jumps.}

\step
\recolor{lift.cell[1][0]}{state=done}
\recolor{lift.cell[1][1]}{state=done}
\recolor{lift.cell[1][2]}{state=done}
\recolor{lift.cell[1][3]}{state=done}
\recolor{lift.cell[1][4]}{state=done}
\recolor{lift.cell[1][5]}{state=done}
\apply{lift.cell[2][0]}{value="2"}
\apply{lift.cell[2][1]}{value="3"}
\apply{lift.cell[2][2]}{value="1"}
\apply{lift.cell[2][3]}{value="5"}
\apply{lift.cell[2][4]}{value="6"}
\apply{lift.cell[2][5]}{value="4"}
\recolor{lift.cell[2][0]}{state=current}
\recolor{lift.cell[2][1]}{state=current}
\recolor{lift.cell[2][2]}{state=current}
\recolor{lift.cell[2][3]}{state=current}
\recolor{lift.cell[2][4]}{state=current}
\recolor{lift.cell[2][5]}{state=current}
\narrate{Row k=2: lift[2][i] = lift[1][lift[1][i]], the node after 4 steps. Since cycle length is 3, jumping 4 steps is the same as jumping 1 step. Node 1 after 4 steps: 1->2->3->1->2, so lift[2][1]=2. Notice it matches lift[0][1] -- that is 4 mod 3 = 1 step.}

\step
\recolor{lift.cell[2][0]}{state=done}
\recolor{lift.cell[2][1]}{state=done}
\recolor{lift.cell[2][2]}{state=done}
\recolor{lift.cell[2][3]}{state=done}
\recolor{lift.cell[2][4]}{state=done}
\recolor{lift.cell[2][5]}{state=done}
\narrate{Binary lifting table complete. We can now jump any node forward by any distance d in O(log d) time by decomposing d into powers of 2. Now let us answer queries.}

\step
\recolor{g.node[1]}{state=current}
\highlight{g.node[3]}
\apply{ans.cell[0]}{value="?"}
\recolor{ans.cell[0]}{state=current}
\narrate{Query 1: distance from node 1 to node 3? Both are in component 1 (same cycle). We need to follow edges from 1 and count steps until we reach 3.}

\step
\recolor{g.node[1]}{state=current}
\recolor{g.edge[(1,2)]}{state=current}
\recolor{g.node[2]}{state=done}
\highlight{g.node[3]}
\narrate{Step 1: teleport from 1 to 2. We have not reached 3 yet. Continue.}

\step
\recolor{g.node[1]}{state=done}
\recolor{g.edge[(2,3)]}{state=current}
\recolor{g.node[2]}{state=done}
\recolor{g.node[3]}{state=good}
\recolor{g.edge[(1,2)]}{state=done}
\narrate{Step 2: teleport from 2 to 3. We reached the target! Distance = 2. On the cycle of length 3, the forward distance from position 0 to position 2 is (2 - 0) mod 3 = 2.}

\step
\apply{ans.cell[0]}{value="2"}
\recolor{ans.cell[0]}{state=good}
\narrate{Answer for query (1, 3) = 2. Two teleportations: 1 -> 2 -> 3.}

\step
\recolor{g.all}{state=idle}
\recolor{ans.cell[0]}{state=idle}
\apply{ans.cell[0]}{value="?"}
\recolor{g.node[4]}{state=current}
\highlight{g.node[2]}
\recolor{ans.cell[0]}{state=current}
\narrate{Query 2: distance from node 4 to node 2? Node 4 is in component 2, node 2 is in component 1. Are they in the same component?}

\step
\recolor{g.node[4]}{state=error}
\recolor{g.node[2]}{state=error}
\recolor{g.edge[(4,5)]}{state=dim}
\recolor{g.edge[(5,6)]}{state=dim}
\recolor{g.edge[(6,4)]}{state=dim}
\recolor{g.edge[(1,2)]}{state=dim}
\recolor{g.edge[(2,3)]}{state=dim}
\recolor{g.edge[(3,1)]}{state=dim}
\narrate{Different components! Node 4 belongs to cycle {4,5,6} and node 2 belongs to cycle {1,2,3}. No sequence of teleportations can cross between components. Immediately return -1.}

\step
\apply{ans.cell[0]}{value="-1"}
\recolor{ans.cell[0]}{state=error}
\narrate{Answer for query (4, 2) = -1. Unreachable -- different connected components in the functional graph.}

\step
\recolor{g.all}{state=idle}
\recolor{ans.cell[0]}{state=idle}
\apply{ans.cell[0]}{value="?"}
\recolor{g.node[1]}{state=current}
\highlight{g.node[1]}
\recolor{ans.cell[0]}{state=current}
\narrate{Bonus query: distance from node 1 to node 1? Same source and target.}

\step
\apply{ans.cell[0]}{value="0"}
\recolor{ans.cell[0]}{state=good}
\recolor{g.node[1]}{state=good}
\narrate{Answer = 0. No teleportations needed -- you are already there. Note: the cycle distance would be 3 (go all the way around), but the minimum is 0 since a = b.}

\step
\recolor{g.all}{state=idle}
\recolor{ans.cell[0]}{state=idle}
\recolor{lift.cell[0][0]}{state=idle}
\recolor{lift.cell[0][1]}{state=idle}
\recolor{lift.cell[0][2]}{state=idle}
\recolor{lift.cell[0][3]}{state=idle}
\recolor{lift.cell[0][4]}{state=idle}
\recolor{lift.cell[0][5]}{state=idle}
\narrate{To handle tail nodes (trees hanging off cycles), we also need depth information. If b is on a's path toward the cycle, we jump a forward by (depth[a] - depth[b]) steps and check if we land on b. If b is on the cycle, we walk a to the cycle, then compute cyclic distance.}

\step
\recolor{g.node[1]}{state=good}
\recolor{g.node[2]}{state=good}
\recolor{g.node[3]}{state=good}
\recolor{g.node[4]}{state=good}
\recolor{g.node[5]}{state=good}
\recolor{g.node[6]}{state=good}
\recolor{g.edge[(1,2)]}{state=good}
\recolor{g.edge[(2,3)]}{state=good}
\recolor{g.edge[(3,1)]}{state=good}
\recolor{g.edge[(4,5)]}{state=good}
\recolor{g.edge[(5,6)]}{state=good}
\recolor{g.edge[(6,4)]}{state=good}
\narrate{Summary: in a functional graph, binary lifting lets us jump forward efficiently, cycle detection identifies components, and combining both gives O(log n) per query. Total complexity: O((n + q) log n). The key insight is that reachability in a functional graph is strictly forward along the unique path from each node into its cycle.}
\end{animation}
